%-*-latex-*-

\documentclass[a4paper,11pt]{article}

% Language and fonts
%
\usepackage[utf8]{inputenc}
\usepackage[T1]{fontenc}
\usepackage[english]{babel}
\usepackage[hyphens]{url}
\usepackage{hyphenat}
\usepackage{array}
\usepackage{microtype}
\usepackage{ebgaramond}

\title{Tutoring notes}
\date{\oldstylenums{2026}}
\author{Christian Rinderknecht}

\begin{document}

\maketitle

\begin{itemize}

  \item \emph{Hogy vagy?} `How are you?' inform. --- \emph{Hogy van?}
    formal.

  \item \emph{Megvagyok.} `I am all right.' (neutral)

  \item \emph{Rosszul/Beteg vagyok.} `I am ill/sick.' (\emph{rossz}
    `bad')

  \item \emph{ön}: pers. pron. `you' sg. form.; \emph{önök} `you'
    pl. form.

  \item \emph{Mi újság?} `What's up?' inform.

  \item \emph{És veled?} `And what with you?' (\emph{vel-} `with' +
    \emph{-ed} 2nd pers. sg. `you' $\rightarrow$
    \emph{veled}). \emph{velem} `with me', \emph{veled} `with you',
    \emph{vele} `with him/her/it', \emph{velünk} `with us',
    \emph{veletek} `with you [pl.]', \emph{velük} `with them'.

  \item \emph{Semmi különös/extra.} `Nothing special.'; \emph{Minden
  rendben.} `Everything is fine.'

  \item \emph{Milyen idő van?} or \emph{Milyen az idő?} `What's the
    weather?'

  \item \emph{Jól/Rossz idő van.} `The weather is good/bad.'

  \item \emph{\underline{Szeles idő} van} `The weather is windy.' ---
    \emph{szél} `wind'. \emph{Fúj a szél.} `The wind is blowing.'

  \item \emph{Esik az eső.} `The rain is falling.'; \emph{Esik a hó.}
    `The snow is falling.'; \emph{\underline{Esős idő} van.} `The
    weather is rainy.'

  \item \emph{Hideg van.} `It is cold.' or \emph{Hideg idő van.}  `The
    weather is cold.'; \emph{\underline{Felhős idő} van.} `The weather
    is cloudy.'; \emph{\underline{Hűvös idő} van.} `The weather is
    chilly.' \emph{Kicsit beteg vagyok, mert hűvös az idő Egerben.}
    (\emph{mert}: `because')

  \item \emph{Fázom.} `I am/feel cold.' (vb. ind. \emph{fázik} `to be
    cold'). \emph{Meg fáztam.} `I caught a cold.' \emph{Meg vagyok
    fázva.} `I have the flu.'

  \item \emph{Süt a nap.} `The sun shines.'; \emph{\underline{Napos
    idő} van.}  `The weather is sunny.'; \emph{Tiszta az ég.} `Clear
    skies.'

  \item \emph{Hány fok van?} `How many degrees it is?'

\end{itemize}

\newpage

\noindent Example of small talk:

\begin{itemize}

  \item[---] \emph{Szia! Hogy vagy?}

  \item[---] \emph{Szia! Jól vagyok, csak egy kicsit fáradt vagyok. És
  te?}

  \item[---] \emph{Megvagyok, köszönöm. Mi újság?}

  \item[---] \emph{Semmi különös. És veled?}

  \item[---] \emph{Velem sincs semmi extra. Egerben milyen idő van?}

  \item[---] \emph{Itt hűvös az idő. Budapesten hány fok van?}

\end{itemize}

\noindent Resuming the lesson:
\begin{itemize}

  \item \emph{Kérdezhetek [...]?} `May/Can I ask [...]?'
    (\emph{kérdez} vb. `to ask as question' + \emph{-het} potential +
    \emph{-ek} 1st pers. sg. pres. ind.)

  \item \emph{valami} indef. pron. `something' (\emph{vala-}
    `some-/-ever' (undefined existence) + interr. pron. \emph{mi}
    `what')

  \item \emph{Mit jelent az, hogy [...]?}  `What does [...] mean?'
    (\emph{jelenteni} vb. inf. `signify', `mean', `report'
    $\rightarrow$ \emph{jelent} 3rd
    pers. sg. ind. pres. indef.). \emph{Ez jelent valamit.} `This
    means something.' \emph{Ez sok\textbf{at} jelent.} `This means a
    lot.'

  \item \emph{mondják} `they say/tell' or pass. `is said/told'
    (vb. \emph{mond} `to say/tell' + \emph{-ják} 3rd
    pers. pl. ind. pres. def.) \emph{Hogy mondják [...]?} `How to say
    [...]?'

  \item \emph{Világos minden?} `Is everything/all clear/understood?';
    \emph{világ} `light (arch.)' (also `world') + \emph{-os} noun
    adj.-form. `-ish' (ornative: `endowed with') $\rightarrow$
    `bright', `understood'.

  \item \emph{fogalmam} `my idea'; \emph{fogalom} `concept, idea'+
    \emph{-am} 1st pers sg. single\hyp{}poss. `my'.

  \item \emph{sincs} `not ... either', `doesn't have either', `there
    isn't either' (\emph{sem} + \emph{nincs} $\rightarrow$
    \emph{sincs}). Use 3rd pers.

    \item \emph{Elnézést!} `Excuse-me!', `Sorry!'. \emph{el-}
      vb. pref. `away, past, over' + \emph{néz} vb. `to look'
      $\rightarrow$ `to overlook or forgive a fault'.
\end{itemize}

Some useful expressions:
\begin{itemize}

  \item \emph{Meg vagyok fázvan.} `I have the flu.'; \emph{Megfáztam.}
    `I got a cold.'

  \item \emph{A számlát kérem.} `The bill, please.' (definite);
    \emph{A kávét kérek.} `I would like a coffee.' (indefinite);
    \emph{Szeretnék a kávét.} (polite)

  \item \emph{Nekem is kell egy tea.} `I also need a tea.' (no
    accusative)

  \item \emph{Nem annyi\textbf{ra} hideg.} `No so cold.'

  \item \emph{Menn\textbf{em} kell.} `I have to go.' (\emph{menni}
    vb. inf. `to go': \emph{megyek}, \emph{mész}, \emph{megy},
    \emph{megyünk}, \emph{mentek}, \emph{mennek})

  \item \emph{Tele vagyok.} `I am full.'; \emph{Elég vagyok.} `I had
    enough.'

  \item \emph{Semmi baj.} `No harm done.' (After an apology)

\end{itemize}

\noindent Countries, nationalities and languages.
\begin{enumerate}

  \item \textbf{Countries ending in \emph{-ország}}: nationality and
    language without ending, e.g. \emph{Franciaország},
    \emph{francia}, \emph{francia}. When an adverb, the language takes
    the suffixes \emph{-ul}/\emph{-ül}, depending on vowel harmony,
    e.g. \emph{\underline{Kicsit beszélek} franci\textbf{ául}.} `I
    speak a bit of French'. (\emph{Kicsit beszélek} forms a group.)

  \item \textbf{Countries ending in \emph{-ia}}: nationality and
    language without ending, e.g. \emph{Dánia} (Danemark), \emph{dán}
    (Danish citizen), \emph{dán} (Danish language). Exceptions:
    \emph{Bulgária}, \emph{bolgár}, \emph{bolgár}; \emph{Ausztria},
    \emph{osztrák}.

  \item \textbf{Nationalities and languages ending in \emph{-i}}:
    \emph{Amerika}, \emph{amerika\emph{i}}; \emph{Kína},
    \emph{kína\emph{i}}.

  \item \textbf{The rest}: \emph{Belgium}, \emph{belga}; \emph{Japán},
    \emph{japán}, \emph{japán}; \emph{Anglia}, \emph{angol},
    \emph{angol}.

\end{enumerate}

\noindent Use \emph{van} when referring to locations or emotions that
are adverbs, e.g. \emph{Ő korházban \textbf{van}?} `Is (s)he at the
hospital?'; \emph{Ő jól \textbf{van}?} `Is (s)he well?'

\begin{itemize}

  \item \emph{mérnök} `engineer'

\end{itemize}

\end{document}
